\documentclass[12pt]{article}
\usepackage[utf8]{inputenc}
\usepackage[T1]{fontenc}
\usepackage{hyperref}
\title{Tahaggart Tamahaq Language Resources for K-12 Curriculum Development}
\author{Compiled by Claude Sonnet 5 (claude-sonnet-5) for Andr\'as Kornai}
\date{}
\begin{document}
\maketitle

\section{Introduction}
Tahaggart Tamahaq (ISO 639-3 thv) is a Tuareg (Berber) language of southern Algeria, spoken mainly in the Ahaggar (Hoggar) area. Sources give conflicting speaker counts, from 40,000 in Algeria to 130,000. Searches found a Tamahaq--French dictionary project and scripture material, but no school curriculum.

\section{Primary Educational Resources}
No primers or textbooks turned up in the searches.

\section{Linguistic and Reference Materials}
\begin{itemize}
\item Tamahaq online dictionary on Webonary. It is an unfinished Tamahaq--French dictionary for the Hoggar variety, built on the work of de Foucauld, Cortade, Mammeri and Barr\`ere: \url{https://www.webonary.work/tamahaq/overview/introduction/?lang=en}
\item General overview: \url{https://en.wikipedia.org/wiki/Tamahaq_language}
\end{itemize}

\section{Religious and Community Resources}
Scripture Earth lists Bible material in Tamahaq (THV): \url{https://www.scriptureearth.org/00eng.php?iso=thv}

\section{Digital and Audio Resources}
\begin{itemize}
\item Global Recordings Network entry (the page title refers to Ghat): \url{https://www.globalrecordings.net/en/language/17129}
\item Endangered Languages Project entry: \url{https://endangeredlanguages.com/node/94394}
\end{itemize}

\section{Curriculum Development Context}
The existing material consists of a historical-lexicon-based dictionary and scripture. A curriculum would need to work out the script (Latin, Tifinagh or Arabic) and the dialect, since Tahaggart, Ajjer and Ghat vary. It would also need to confirm how many speakers there are.

\end{document}